% 本文件由 LaTeX 排版 Agent 根据有效 translation.json 自动生成。
% author=Justin Martyr; work_id=0135; title=致希腊人论辞

\chapter{致希腊人论辞}

% structure: p0001 source=h1 target=omitted reason=page-title-already-rendered-by-parent

% structure: p0002 source=h2 target=section reason=relative-heading-rank; base=h2

\section{第一章 \Person{尤斯定}说明他离开希腊习俗的理由}

希腊人啊，不要以为我脱离你们的习俗是毫无理由和轻率的；因为我在其中找不到任何神圣或蒙天主悦纳的事。你们诗人的作品本身就是疯狂和放纵的纪念碑。因为任何成为你们最杰出导师的学生，都比其他人更深陷困境。首先，他们说\Person{阿伽门侬}为了迎合他兄弟的过度情欲、疯狂和不受约束的欲望，甚至随意牺牲了自己的女儿，并为了解救被麻风牧羊人劫走的\Person{海伦}而搅动整个希腊。但在战争期间抓获俘虏时，\Person{阿伽门侬}自己却被\Person{克律塞伊斯}俘虏，并为了\Person{布里塞伊斯}而与\Person{忒提斯}之子结怨。而\Person{佩里德斯}本人，他渡过河流，推翻了\Place{特洛伊}，征服了\Person{赫克托耳}——你们这位英雄却成了\Person{波吕克塞娜}的奴隶，被一个死去的\Place{亚马孙}女战士征服；他脱下神造的铠甲，穿上婚袍，在\Place{阿波罗神庙}中成为爱情的祭品。而\Place{伊萨卡}的\Person{尤利西斯}则把恶习美化为美德。的确，他驶过\Person{塞壬}海妖身旁，证明他缺乏值得称道的谨慎，因为他无法依靠自己的谨慎来堵住耳朵。\Person{特拉蒙}之子\Person{埃阿斯}，手持七层牛皮盾牌，在与\Person{尤利西斯}争夺盔甲的竞赛中失败后发了疯。我可不想受这样的教导。我不贪求这样的美德，以致相信\Person{荷马}的神话。因为整部史诗——\WorkTitle{伊利亚特}和\WorkTitle{奥德赛}的开头和结尾——都是一个女人。\par

% structure: p0004 source=h2 target=section reason=relative-heading-rank; base=h2

\section{第二章 揭穿希腊神谱}

但继\Person{荷马}之后，\Person{赫西奥德}写了他的\WorkTitle{工作与时日}，谁会相信他那胡言乱语的神谱呢？因为他们说，\Person{克洛诺斯}，\Person{乌拉诺斯}之子，起初杀死了他的父亲，并夺取了他的统治权；又因害怕自己遭受同样的命运，便宁愿吞食自己的孩子；但靠着\Person{枯瑞忒斯}的计谋，\Person{朱庇特}被偷偷带走并隐藏起来，后来用锁链捆绑了他的父亲，并瓜分了帝国：据说\Person{朱庇特}得到了天空，\Person{尼普顿}得到了深海，\Person{普路托}得到了\Place{哈得斯}的部分。但\Person{普路托}抢走了\Person{普洛塞庇娜}；\Person{刻瑞斯}在荒野中流浪寻找她的孩子。这个神话在\Place{厄琉息斯}圣火中被传颂。同样，\Person{尼普顿}在\Person{墨拉尼珀}打水时强占了她，此外还侮辱了无数海中宁芙，若是列举她们的名字，会耗费我们许多话语。至于\Person{朱庇特}，他是个奸夫：化成半兽半人与\Person{安提俄珀}相交，化成黄金与\Person{达娜厄}相交，化成公牛与\Person{欧罗巴}相交；此外，他为了\Person{勒达}还生出了翅膀。对\Person{塞墨勒}的爱既证明了他的不贞，也导致了\Person{塞墨勒}的嫉妒。他们还说他劫走了\Place{弗里吉亚}的\Person{伽倪墨得斯}做他的侍酒者。这就是\Person{萨图尔努斯}之子的功绩。而你们显赫的\Person{勒托}之子（\Person{阿波罗}），自称能预言，却暴露了自己在说谎。他追求\Person{达芙妮}，却没有得到她；对爱他的\Person{许阿铿托斯}，他没有预言他的死亡。我且不提\Person{密涅瓦}的男性特质，也不提\Person{巴克斯}的女性本质，也不提\Person{维纳斯}的淫荡本性。希腊人啊，请把惩治弑亲的法律、通奸的刑罚和鸡奸的耻辱读给\Person{朱庇特}听吧。教会\Person{密涅瓦}和\Person{狄安娜}去做女人的工作，教\Person{巴克斯}去做男人的工作。一个女人披挂盔甲，或一个男人用铙钹、花环、女装装饰自己，身边跟着一群狂女，这有什么体统呢？\par

% structure: p0006 source=h2 target=section reason=relative-heading-rank; base=h2

\section{第三章 希腊神话的愚妄}

至于\Person{赫拉克勒斯}，他因三个夜晚而闻名，被诗人歌颂他成功的功业；他是\Person{朱庇特}之子，杀死狮子，消灭多头\Person{许德拉}；杀死凶猛巨大的野猪，能射杀飞快的食人鸟，从\Place{哈得斯}带出三头犬；彻底清扫了巨大的\Person{奥革阿斯}牛圈，杀死了鼻孔喷火的公牛和母鹿，从树上摘下金苹果，杀死毒蛇（并且出于一个不可言说的原因，杀死了\Person{阿刻罗俄斯}和杀客的\Person{布西里斯}），翻山越岭去取据说能发出清晰言语的泉水；这人有能力完成如此众多、如此伟大、如此惊人的伟业——却多么幼稚地被\Person{萨梯尔}的铙钹震晕，被女人的爱情征服，被大笑的\Person{吕达}击中臀部！最后，由于无法脱下\Person{涅索斯}的衬衣，他自己点燃了自己的火葬堆，就这样死了。愿\Person{武尔坎}收起他的嫉妒，不要因自己年老跛足而遭恨，而\Person{马尔斯}因年轻英俊被爱而感到嫉妒。因此，希腊人啊，你们的神明被判为放纵，你们的英雄沦为女性化，正如你们戏剧所依据的历史所宣告的：\Person{阿特柔斯}的诅咒、\Person{提厄斯忒斯}的床、\Person{珀罗普斯}家的污点、\Person{达那俄斯}因仇恨而杀人并酗酒狂怒使\Person{埃古普托斯}无子，以及由复仇女神摆布的\Person{提厄斯忒斯}盛宴。而\Person{普洛克涅}至今还在飞翔哀鸣；她的\Place{雅典}姐妹被割去舌头发出尖叫声。有什么必要谈论\Person{俄狄浦斯}所受的折磨、\Person{拉伊俄斯}被杀、娶母为妻，以及那些既是他的兄弟又是他的儿子的互相残杀呢？\par

% structure: p0008 source=h2 target=section reason=relative-heading-rank; base=h2

\section{第四章 希腊人的无耻行径}

我也开始憎恶你们的公共集会。那里有纵情的宴饮、挑动情欲的微妙笛声、无益而奢华的涂油，以及花冠加冕。你们以如此众多的罪恶驱逐羞耻；你们以此充塞心智，被放纵所挟持，将邪恶疯狂的淫乱当作常事。我还要对你们说：希腊人啊，当你们的儿子效仿\OriginalTerm{朱庇特}{Jupiter}，起来反抗你们，夺走你们自己的妻子时，你们为何向他发怒？你们为何视他为敌人，却又崇拜与他相似的神明？你们为何谴责妻子不贞，却为\OriginalTerm{维纳斯}{Venus}筑庙供奉？若这些事情是由他人叙述，它们似乎只是诽谤的指控，而非真理。但如今，你们的诗人亲自歌唱这些事，你们的历史也大声宣扬它们。\par

% structure: p0010 source=h2 target=section reason=relative-heading-rank; base=h2

\section{第五章 结语呼吁}

从今以后，希腊人啊，来领受无与伦比的智慧，受教于神圣的圣言，与不朽的君王相识；不要认那些屠杀全民族的人为英雄。因为我们的统治者，神圣的圣言，现今仍不断帮助我们的那一位，并不渴望身体的力量和容貌的美，也不渴求世上贵族的崇高精神，而是一个纯洁的灵魂，为圣洁所坚固，以及我们君王的暗号——神圣的行为，因为藉着圣言，能力进入灵魂。啊，和平的号角，向着争战的灵魂！啊，驱散可怕情欲的武器！啊，熄灭灵魂内在之火的训诫！圣言所施加的影响，并不造就诗人：它不装备哲学家或娴熟的演说家，而是藉其训导使必死的成为不朽，使必死的成为神；并从大地将他们运送到奥林波斯之上之境。来，受教；成为像我一样的，因为我以前也像你们一样。这些征服了我——训导的神性和圣言的能力：正如一位娴熟的弄蛇者将可怕的爬虫从洞穴中诱出并使其逃遁，同样，圣言将我们感性本性的可怕情欲从灵魂的最深处驱赶出去；首先驱走私欲，由此生出一切邪恶——仇恨、争斗、嫉妒、争竞、愤怒，以及诸如此类。私欲一旦被逐出，灵魂便变得平静安宁。它从那淹没至颈的邪恶中得释放，就回到造它的那一位那里。因为它理应恢复到它起初离开的状态，一切灵魂从那里出发或仍在其中。\par

\SourceRecord{Translated by Marcus Dods.；Ante-Nicene Fathers；Vol. 1.；Edited by Alexander Roberts, James Donaldson, and A. Cleveland Coxe.；Buffalo, NY: Christian Literature Publishing Co.,；1885.；<http://www.newadvent.org/fathers/0135.htm>.}
